\begin{abstract}
배틀로얄 게임의 공개 텔레메트리는 줄어드는 활동 영역에서 다수의 팀이 경쟁하는 과정을 대규모 시공간 데이터로 기록한다.
본 연구는 단계 경계에서 생존한 팀이 다음 단계 안에 탈락할 조건부 확률을 서식지 축소 관점의 단계 조건부 이산시간 hazard로 정식화하고, 이 위험의 예측을 무엇이 좌우하는지 분석한다.
PUBG 13,948경기의 텔레메트리 약 3억 2천만 건에서 예측 시점 이후의 정보가 섞이지 않는 팀--단계 표본 817,766개를 구성하고, 12개 코호트에서 \rev{두 기존 설계의 적용판과 같은 데이터·평가 조건으로} 비교한 뒤 설계 요소를 하나씩 바꾸어 차이를 분해하였다.
\rev{주변 경쟁 밀도와 이동 부담이 클수록, 남은 팀원이 적을수록 위험이 높았고, 구역 안쪽 깊이와 위험의 음의 연관은 구역이 작아질수록 강해졌다.
간이 모형의 네 변수는 같은 단계 안의 비교에서도 고차원 모델 판별력의 약 71\%를 재현하였다.}
\rev{본 연구의 설계는 모든 코호트에서 두 적용판보다 AUC가 높았고(0.778--0.779 대 0.686--0.763) 같은 단계 안의 비교에서도 앞섰다.
상대위험 기반 적용판은 단계별 위험 수준을 거의 평평하게 예측했지만(단계별 보정 오차 0.064, 기저위험 결합 시 0.017, 본 설계 0.013), 원래 과제인 우승 팀 예측에서는 가장 높아 학습 목적이 추정 대상을 정함을 보였다.
설계 간 차이는 주로 단계별 기저위험·학습 목적과 서식지 맥락이 설명하였다.
모델 구조는 같은 경기 안의 순위를 바꾸지 않았고, 같은 단계에서 경기마다 다른 위험 수준을 담는 데 작게(단계 안 AUC 약 0.01) 기여하였다.}
\end{abstract}

% 주제어는 EasyChair 제출 양식에 입력(본문에서는 분량 때문에 생략):
% 시공간 텔레메트리, 이산시간 hazard, 데이터 누수, 생존분석, e스포츠 분석
